% CoRL 2026 Workshop on Human-Centered Robot Learning and Interaction
% Extended abstract, 4 pages excluding references. Non-archival, single-blind.
%
% TODO before submitting:
%   1. Fill in the surname on line \author{}.
%   2. Decide with your advisor whether he is a co-author on this.
%   3. Check the workshop's OpenReview page for a required style file; if one
%      exists, move this content into it rather than using article class.

\documentclass[10pt]{article}

\usepackage[letterpaper,margin=1in]{geometry}
\usepackage{amsmath,amssymb}
\usepackage[round]{natbib}
\usepackage{booktabs}
\usepackage{microtype}
\usepackage[hidelinks]{hyperref}
\usepackage{titlesec}
\usepackage{caption}

\captionsetup{font=small,labelfont=bf,skip=6pt}
\titlespacing*{\section}{0pt}{10pt}{4pt}
\titlespacing*{\subsection}{0pt}{8pt}{3pt}
\setlength{\parskip}{2pt}

\title{\vspace{-2.2em}Decision-Relevant Preference Elicitation\\
under a Population Prior}

\author{Joshua [SURNAME]\\
\small University of Southern California\\
\small \texttt{jjt\_373@usc.edu}}

\date{}

\begin{document}
\maketitle
\vspace{-2.2em}

\begin{abstract}
\noindent
Active preference-based reward learning selects the comparison that reveals the
most about a user's reward function. We argue this is the wrong target: much of
what a reward function encodes never changes what the robot does, so the robot
spends a person's limited patience resolving uncertainty with no behavioral
consequence. We select queries instead by expected reduction in policy regret,
computed under a prior learned from a population of previous users. In an MDP
where optimal policies and true regret are exact, decision-relevant selection
reaches lower regret than information-gain selection after ten queries
($2.37 \pm 0.26$ vs.\ $3.39 \pm 0.33$, $n{=}150$), and the paired advantage in
queries-to-threshold is $0.63$ queries, 95\% CI $[0.02, 1.25]$ --- positive but
fragile. We also report two negative results and one retraction. The idea is
not new; regret-based reward elicitation dates to 2009. What is new is the
instantiation with human trajectory comparisons and a learned population prior,
and we are explicit about how thin the current evidence for it is.
\end{abstract}

\section{Introduction}

A robot can learn what a person wants by showing them pairs of behaviors and
asking which is better. The bottleneck is the person: physical robots cannot
afford the thousands of comparisons a simulated agent can~\citep{sadigh2017}.
Nearly all work on this bottleneck asks the same question --- which comparison
is most \emph{informative about the reward function}~\citep{biyik2018,biyik2019,
biyik2024,hejna2022}.

We think that target is misspecified. A reward function's purpose is to induce a
policy, and the map from rewards to policies is heavily many-to-one. Two reward
estimates that differ substantially often produce identical behavior. Every
query spent separating them buys nothing the robot will ever act on.

The alternative is to select queries by how much the answer changes what the
robot would \emph{do}. This is not a new idea. \citet{regan2009} posed exactly
this problem for MDPs seventeen years ago, minimizing the precision with which a
reward must be specified while still producing near-optimal policies, and
selecting queries by minimax regret reduction. What has not been done is the
instantiation that matters for robot learning: human trajectory-segment
comparisons rather than bound queries on reward values, and a \emph{prior
learned from previous users} rather than an uninformative prior over all of
reward space. The population prior is what makes the approach affordable ---
we search only the region real users occupy.

\textbf{Contributions.} (i) A value-of-information acquisition function over
trajectory comparisons that scores queries by expected reduction in policy
regret under a population prior. (ii) An evaluation in an MDP where optimal
policies and regret are computed exactly, so that the effect under test is not
confounded with approximation error. (iii) An honest account of what we can and
cannot yet show, including two negative results and one retracted claim.

\section{Related work}

\textbf{Decision-relevant elicitation already exists.} \citet{regan2009}
introduced regret-based reward elicitation for MDPs; \citet{alizadeh2015}
extended it approximately. \citet{modelfreepe2024} score queries by expected
value of information toward recommendation quality, in recommender systems with
no policy. We claim no conceptual novelty over this line.

\textbf{Acquisition functions in robot reward learning} form a sequence of
corrections. Volume removal~\citep{sadigh2017} admits the trivial query --- two
identical options --- as a global optimum. Information gain~\citep{biyik2019}
makes it a global minimum and, by subtracting the human's expected answer
uncertainty, yields questions people can actually answer. \citet{biyik2024}
argue parameter-space entropy is wasteful because many parameters give one
reward and many rewards give one behavior; \citet{bickfordsmith2023} make the
same argument for active learning generally, showing information gain about
parameters can underperform random selection. Our objective is the natural
endpoint of both critiques, moved one step further: not to predictions, but to
actions.

\textbf{Population priors.} Latent-user models learn how users
differ~\citep{poddar2024,pal2024,maxmin2024}, but adapt from fixed or randomly
chosen queries. We reuse their premise --- that a population constrains who a
new user can be --- and spend it on choosing questions.

\section{Method}

\textbf{Setup.} Reward is linear in state features, $r(s) = w \cdot \phi(s)$,
with $\|w\| = 1$. A trajectory segment is summarized by its discounted feature
count $f = \sum_t \gamma^t \phi(s_t)$; a query is a pair of segments with
$\delta = f_A - f_B$. A user with consistency $b > 0$ answers according to
Bradley--Terry, $P(A \succ B) = \sigma(b \, w \cdot \delta)$.

Normalizing $w$ is not cosmetic. Only the product $b\|w\|$ affects any answer
probability, so leaving both free makes consistency an artifact of the
parameterization rather than a property of the person.

\textbf{Population prior.} We represent beliefs about the current user as
weighted particles $\{w_i\}_{i=1}^N$ drawn from a multi-modal prior fit to
previous users, updated by exact Bayes on each answer. This is the component
absent from \citet{regan2009}: the elicitor never searches reward space at
large, only the part the population occupies.

\textbf{Baseline acquisition (BALD).} Information about the reward,
\begin{equation}
\text{score}(q) = H\!\left[\mathbb{E}_w\, p(y \mid w)\right]
                - \mathbb{E}_w H\!\left[p(y \mid w)\right],
\end{equation}
high when particles disagree, with per-particle answer noise subtracted. The
subtraction is what makes the trivial query score exactly zero.

\textbf{Decision-relevant acquisition (VOI).} Let $\pi(\rho)$ be the optimal
policy for the mean reward of posterior $\rho$, and define the loss of acting on
it as the expected regret across who the user might be,
\begin{equation}
L(\rho) = \mathbb{E}_{w \sim \rho}\!\left[V^{*}_{w} - V^{\pi(\rho)}_{w}\right].
\end{equation}
We then score each candidate by the expected improvement in that loss,
\begin{equation}
\text{score}(q) = L(\rho) - \sum_{y} P(y)\, L(\rho \mid y).
\end{equation}
Each candidate costs two value iterations. That expense is why the 2009 work
remained in small tabular MDPs, and why the population prior matters here.

\section{Experiments}

\textbf{Domain.} A $6{\times}6$ gridworld with one-hot terrain features,
$\gamma = 0.95$, in which the optimal policy and the true regret of acting on a
wrong reward are both exact. Approximation error in a larger domain would be
indistinguishable from the effect we are measuring, so the first experiment
admits none. Users are drawn from a three-mode population prior, which is also
the elicitor's particle prior.

\textbf{Protocol.} Four strategies share the same posterior machinery, the same
users and the same candidate queries: \textsc{random}; \textsc{bald};
\textsc{voi}; and an \textsc{oracle} that cheats by scoring against true regret.
The oracle is essential --- without it one cannot distinguish a good method from
an easy problem. Runs are paired, so we report paired bootstrap confidence
intervals (10{,}000 resamples) rather than independent standard errors.

\begin{table}[t]
\centering
\small
\begin{tabular}{lcccc}
\toprule
& \textsc{random} & \textsc{bald} & \textsc{voi} & \textsc{oracle} \\
\midrule
Regret after 10 queries & $3.67 \pm 0.34$ & $3.39 \pm 0.33$ & $\mathbf{2.37 \pm 0.26}$ & $0.25 \pm 0.07$ \\
Queries to threshold    & $6.08 \pm 0.34$ & $5.77 \pm 0.33$ & $\mathbf{5.14 \pm 0.33}$ & $2.21 \pm 0.21$ \\
Runs reaching threshold & 66\% & 69\% & 75\% & 95\% \\
\bottomrule
\end{tabular}
\caption{Results at $d{=}8$ over 150 user-runs (10 seeds $\times$ 15 users),
mean $\pm$ standard error. Lower is better.}
\label{tab:main}
\end{table}

\textbf{Results.} Table~\ref{tab:main} gives the powered run. On the paired
comparison of queries-to-threshold, \textsc{voi} improves on \textsc{bald} by
$0.63$ queries, 95\% CI $[0.02, 1.25]$; on \textsc{random} by $0.94$, CI
$[0.18, 1.71]$. \textsc{bald} does \emph{not} separate from \textsc{random} in
this domain, CI $[-0.43, 1.03]$.

We do not want to oversell the first of those. The interval's lower bound is
$0.02$, $43\%$ of runs tie, and a third of runs never reach the threshold, so
censoring may bias the estimate. The final-regret comparison is cleaner --- a
roughly three-standard-error gap --- and we expect it to become the primary
metric in future runs, a decision we will fix before collecting them rather
than after.

\textbf{The oracle gap is the more useful number.} \textsc{oracle} reaches the
threshold in $2.21$ queries against \textsc{voi}'s $5.14$, and ends at regret
$0.25$ against $2.37$. Our acquisition captures roughly half of what is
available even with perfect knowledge of the user. The remaining headroom is
about estimator quality --- we currently use only the posterior-mean policy and
a two-point marginalization over the answer --- rather than about the choice of
target.

\section{What did not work}

\textbf{At $d{=}4$, the effect is not detectable.} Over 48 user-runs the
\textsc{bald}$-$\textsc{voi} difference was $0.42$ queries, CI $[-0.38, 1.17]$.
Sweeping feature dimension gives $+0.42$ at $d{=}4$, $+1.42$ at $d{=}8$ and
$+1.00$ at $d{=}12$ at 24 runs per cell, with every interval containing zero.
The direction is consistent with reward and behavioral uncertainty decoupling as
dimension grows, but we cannot presently claim it.

\textbf{A related factorization gave nothing.} In separate experiments we split
the user latent into preference direction and answer consistency, expecting the
second to matter when population reliability varies. Against a single-latent
baseline it did not improve held-out prediction at any spread, and recovered
consistency at correlation near zero. A follow-up showed why: consistency is
recoverable ($r = 0.87$) when the preference direction is supplied and query
difficulty is controlled, but easy queries are catastrophic for it (median
absolute error $23$ against a true median of $3$), because a user who answers
every easy question correctly is indistinguishable from a perfect one. The
bottleneck is joint estimation from a small uncontrolled context set.

\textbf{A retraction.} Our first single-seed run showed \textsc{bald}
performing worse than \textsc{random}, which would have echoed
\citet{bickfordsmith2023} neatly. It did not survive additional seeds. We
report it because the workshop invites negative results, and because the
temptation to keep a convenient single-seed finding is the failure mode this
venue exists to discourage.

\section{Limitations and next steps}

Everything here is simulation with users generated from the same response model
the estimator assumes --- the most generous possible condition, and one that
never holds with people. The domain is small enough for exact regret, which is
also small enough that reward and behavioral uncertainty may largely coincide.
Immediate next steps are a better VOI estimator, a domain where the decoupling
is structural rather than incidental, and a mismatch protocol that samples
answers from response models outside the assumed family.

The question we would most like feedback on: given that the decision-relevant
criterion is seventeen years old, is the population-prior instantiation a
sufficient contribution on its own, or does it need accompanying theory?


\begin{thebibliography}{99}
\small
\bibitem[Alizadeh et al.(2015)]{alizadeh2015} P.~Alizadeh et al.
Approximate regret based elicitation in Markov decision processes. \emph{RIVF}, 2015.

\bibitem[Bickford Smith et al.(2023)]{bickfordsmith2023} F.~Bickford Smith et al.
Prediction-oriented Bayesian active learning. \emph{AISTATS}, 2023.

\bibitem[B\i y\i k and Sadigh(2018)]{biyik2018} E.~B\i y\i k and D.~Sadigh.
Batch active preference-based learning of reward functions. \emph{CoRL}, 2018.

\bibitem[B\i y\i k et al.(2019)]{biyik2019} E.~B\i y\i k et al.
Asking easy questions: a user-friendly approach to active reward learning.
\emph{CoRL}, 2019.

\bibitem[B\i y\i k et al.(2024)]{biyik2024} E.~B\i y\i k et al.
A generalized acquisition function for preference-based reward learning.
\emph{ICRA}, 2024.

\bibitem[Hejna and Sadigh(2022)]{hejna2022} J.~Hejna and D.~Sadigh.
Few-shot preference learning for human-in-the-loop RL. \emph{CoRL}, 2022.

\bibitem[Chaudhari et al.(2024)]{modelfreepe2024} S.~Chaudhari et al.
Model-free preference elicitation. \emph{IJCAI}, 2024.

\bibitem[Chakraborty et al.(2024)]{maxmin2024} S.~Chakraborty et al.
MaxMin-RLHF: alignment with diverse human preferences. \emph{ICML}, 2024.

\bibitem[Poddar et al.(2024)]{poddar2024} S.~Poddar et al.
Personalizing RLHF with variational preference learning. \emph{NeurIPS}, 2024.

\bibitem[Chen et al.(2024)]{pal2024} D.~Chen et al.
PAL: sample-efficient personalized reward modeling for pluralistic alignment,
2024.

\bibitem[Regan and Boutilier(2009)]{regan2009} K.~Regan and C.~Boutilier.
Regret-based reward elicitation for Markov decision processes. \emph{UAI}, 2009.

\bibitem[Sadigh et al.(2017)]{sadigh2017} D.~Sadigh et al.
Active preference-based learning of reward functions. \emph{RSS}, 2017.
\end{thebibliography}

\end{document}
